\paragraph{Claims-as-builds.}
This book adopts a claims-as-builds workflow: each substantive claim is paired with an executable artifact, test, or check that can be run, inspected, or reproduced. In practice, a theorem is not treated as complete until it is accompanied by the corresponding proof object, derivation script, verification harness, or other machine-checkable witness appropriate to the topic at hand. This convention keeps the exposition tightly coupled to reproducible evidence and makes the logical dependencies between results explicit.

\paragraph{Acquire--Calibrate--Log.}
The working protocol for examples, experiments, and derived results has three stages. First, acquire the relevant specifications, datasets, or source artifacts. Second, calibrate the chapter-specific error model \(\varepsilon(r)\) at the appropriate resolution \(r\), so that approximations, tolerances, and comparison criteria are stated consistently across the book. Third, log reproducibility metadata, including hashes, random seeds, toolchains, and other build parameters needed to reconstruct the result. This protocol is used uniformly for analytic derivations, computational experiments, and build logs.

\paragraph{Reading map.}
The book is organized as a dependency graph rather than a strictly linear narrative, so readers may enter from different backgrounds and follow different paths through the material. The reading map indicates which chapters depend on which earlier constructions, allowing a reader to move from foundations to applications, or to jump directly to a topic of interest while consulting prerequisite material as needed. The map is meant to support both sequential reading and selective reference use, with the dependency structure serving as a guide to conceptual prerequisites and technical reuse.

\paragraph{Usage convention.}
Throughout the book, when a result is stated, the accompanying artifact should be understood as part of the claim. When a result is approximate, the relevant calibration and error bookkeeping should be stated explicitly. When a result is reproduced, the logged metadata should be sufficient to recover the same computational environment and verify the same outputs. Readers should therefore expect proofs, code, data, and build records to function together as the evidential basis for the text.

\paragraph{How to read this book.}
For readers with a mathematical background, the recommended path is to begin with the foundational chapters, then follow the dependency graph into modularity, feedback, and symmetry. For readers coming from engineering or implementation practice, it is often useful to start with the usage conventions here, consult the notation chapter for symbols and error metrics, and then move into the chapters whose artifacts are most directly relevant to the problem at hand. In either case, the reading map is intended to make prerequisite structure explicit without forcing a single linear route through the material.
